\clearpage
\section{结论}
\label{sec:conclusion}

本实验以 GCC/Clang 编译工具链、LLVM IR、RISC-V 汇编及 AscendNPU IR 为研究对象，通过编译产物分析、程序运行验证和中间表示对照，研究了源程序到目标程序的转换过程及不同抽象层次下的程序语义表示。

在完整编译流程的研究中，通过对预处理文件、LLVM IR、汇编程序和 ELF 目标文件的分析，明确了预处理、编译、汇编与链接各阶段的功能及其相互关系。不同优化配置下的实验进一步表明，编译优化不仅改变指令数量，也会重构程序的控制流和数据流。Clang 在 O2 配置下将原有循环识别为整数归约，并进行向量化、循环展开和标量尾部处理；GCC 与 Clang 在相同优化级别下则表现出不同的代码生成结果。受控对照进一步发现，向量化与展开的结构影响不能简单相加；GCC 在本例中显式启用向量化或使用 O3 仍未生成向量循环，而 Clang 的 noinline 变体因循环内保留的辅助调用未能向量化。这说明编译结果受到优化策略、编译器实现及目标架构约束的共同影响，不能仅依据优化级别或代码大小评价编译效果。

在 SysY 等价实现实验中，针对两个覆盖整数、浮点、控制流、函数和数组等语言特性的示例程序，分别构建了源语言、LLVM IR 和 RISC-V 汇编三种实现。三种实现通过相同输入及独立预期进行验证，共完成 480 次运行测试，结果全部符合预期。结合控制流图、SSA 结构和 RISC-V 调用约定的分析可以发现，源语言的语义一致性不仅依赖算术结果相同，还依赖控制流分支、短路求值的副作用、变量状态合并、数据布局和函数调用约定的正确实现。针对短路求值、有符号余数和浮点转换构造的错误变体均被测试检出，进一步验证了测试方法对这些特定语义错误的辨识能力。

在进阶实验中，通过对 AscendNPU IR VecAdd 示例的逐层 lowering 分析，观察了多级中间表示如何在保持向量加法计算语义的基础上，逐步引入设备执行所需的约束。核类型推导确定了 AIV 执行目标，内存规划将逻辑缓冲映射至 UB 地址空间，并通过满足原地执行条件的存储复用减少静态缓冲占用；跨流水同步和模板转换则进一步明确了数据搬运、向量计算及设备操作之间的依赖关系。这表明，多级中间表示不仅用于表达计算本身，也能够为面向特定硬件的存储管理和执行调度提供逐步细化的表示基础。

本实验的结论仍受到工具版本、测试输入和执行环境的限制。SysY 等价性验证覆盖的是指定有效输入域内的有限样本，不能替代完整的形式化等价证明；QEMU 执行结果不代表真实 RISC-V 硬件性能；AscendNPU IR 实验完成了主机侧中端变换，但尚未生成可执行的设备目标文件，因此不能据此判断设备端数值正确性或性能。总体而言，本实验建立了从源程序、LLVM IR、目标汇编到设备相关中间表示的分析联系，并通过可追溯的编译产物和运行结果，对不同编译阶段的功能、语义保持机制及中端 lowering 过程进行了验证。
